% Ad Familiares 11.7 --- headnote
% ad-familiares-11-07, late 44 / early 43 BC, Rome

\textit{Cicero to Decimus Brutus, consul-designate and \textit{imperator}, written from Rome --- Perseus dateline \textit{Scr. Romae xiii K. \dots a. 710 (44)} (the month in Perseus's transmission is corrupted; the entry in \texttt{meta/works.yaml} carries a year-precision placeholder of 16 April 44 BC, but the letter's reference to the Martian and Fourth Legions having already declared Antony a public enemy puts it after their defection in November 44, and most editors place it in late December 44 or early January 43, before D. Brutus had taken the field). D. Brutus, one of the principal Liberators of the Ides, was now besieged at Mutina by Antony and waiting on signals from Rome about how aggressively to press the war.}

\textit{Cicero's argument is the political brief he had spent the autumn putting together: the Senate is not yet free, so do not wait for its formal authority; the act that frees the Republic is not authorised in advance, it is ratified by what it accomplishes. The Ides of March were the first such act, the raising of the new army the second, and Brutus must now follow through without flinching. The four parallel clauses of section 2 --- that to hesitate now would condemn his own act, would condemn Octavian (``the young man, or rather the boy, Caesar'') for taking up the public cause on private judgement, and would condemn the veteran soldiers and the two legions for the same --- is Cicero's tight-knit defence of the whole irregular coalition that has come together against Antony. The closing maxim, that the will of the Senate ought to be taken for its authority when its authority is hindered by fear, is the constitutional formula on which the campaign rests.}
